\section{Introduction}

Automatic presentation generation aims to turn 
%source documents, papers, or 
user requests into structured slide decks, and 
has gained increasing attention because slides are widely used, yet creating high-quality presentations remains time-consuming and cognitively demanding
\citep{sun2021d2s,maheshwari2024presentations,mondal2024personaaware,zeng2025slidetailor}.
Recent agentic systems further advance this task by 
%planning, constructing, evaluating, and refining 
producing 
complete decks through multi-modal or tool-based workflows
\citep{ge2025autopresent,zheng2025pptagent,xu2025pregenie,xie2025slidebot,liang2025slidegen,zheng2026deeppresenter,ozden2026arcdeck}.
Although these systems can now produce complete and visually polished decks,
they still lack persistent personalization, which is essential for generating ready-to-use slide decks because users vary in domain, purpose, style, and presentation habits.
For example, users may prefer different layouts, templates, and styles when creating slides for different purposes, such as academic presentations versus business presentations.
%An effective personalized slides generation framework should build and maintain user profiles, learning how a particular user tends to organize, style, and revise presentations for different intents from interaction history, rather than relying on users' repeatedly prompting their preferences. 
An effective personalized slide generation framework should build and maintain user profiles that capture users' long-term preferences in organizing, styling, and revising presentations across different intents, rather than requiring users to repeatedly specify their preferences in every interaction.


%: they do not reliably learn how a
%particular user tends to organize, style, and revise presentations across
%interactions. 

%building and updating user profiles from interaction history, especially revision
%feedback, rather than relying only on a single prompt or task-time template.


Prior work has progressively expanded the capability of presentation-generation
systems \citep{sun2021d2s,mondal2024personaaware,
bandyopadhyay2024enhancing,maheshwari2024presentations,ge2025autopresent,
xu2025pregenie,xie2025slidebot,liang2025slidegen,ozden2026arcdeck}. 
PPTAgent moves beyond text-to-slides by coupling generation with
presentation-specific evaluation \citep{zheng2025pptagent}, and DeepPresenter
introduces environment-grounded reflection for agentic presentation generation
\citep{zheng2026deeppresenter}. These systems improve general generation and
agentic refinement, but do not explicitly model user-specific personalization.
SlideTailor addresses personalization by conditioning scientific
slide generation on reference slides and task-time templates
\citep{zeng2025slidetailor}, but its personalization remains tied to provided
examples or template conditions rather than to an accumulated user profile.
These observations therefore expose a central gap: slide generation agents still lack a user-facing multi-turn dialogue process that converts revision requests into reusable preference memory and preserves local revision constraints across later turns.

\begin{figure}[!t]
  \centering
  \includegraphics[width=\linewidth]{pptagent_memory_workflow}
  \vspace{-3mm}
  \caption{Overview of \methodname. The framework comprises long-term memory
  and working memory. Long-term memory stores user profile memory and tool
  memory for persistent personalization and reusable execution experience,
  while working memory carries the current session state for personalized
  generation and localized revision. In round $t$, $s_t$ is the current
  session state, $f_t$ is the user feedback, and $s_{t+1}$ is the
  updated state after \emph{Modify Exec}, the memory-guided revision executor.}
  \label{fig:pptagent_memory_workflow}
\end{figure}

This gap has two roots. First, personalization in presentation authoring is
often revealed through revision rather than fully specified before generation,
yet existing slide generation agents still handle edits by re-contextualizing or
re-generating large parts of the deck. As a result, small changes must compete
with deck state and feedback history for limited context, making multi-turn
local modification fragile. Second, current systems mainly improve
presentation-agent workflows, tools, and evaluators, but still treat
personalization as an implicit byproduct of prompting rather than a direct service enabled by
memory design. In a similar spirit to agent memory work
\citep{zhong2023memorybank,packer2023memgpt,wu2025memoryage}, personalization of 
slide generation can be largely enhanced by incorporating memory as an explicit framework rather than
an undifferentiated dialogue buffer.
%consistent with recent agent-memory work

To address these problems, in this paper, we present MemSlides that introduces scoped slide
locality as its revision strategy for multi-turn slide editing. Rather
than re-reading or re-writing the whole deck for each feedback turn,
\methodname projects the request onto the smallest affected slide region and
operates on a bounded repair surface. It reads only a structured snapshot of
that local surface, including its local layout structure, available selectors,
and exposed style rules, and writes back only patches scoped to explicit
selectors or to those style rules. In this way, both reading and writing stay
local by construction, reducing context pressure and unintended drift while
preserving already aligned content across turns.

Building on this modification strategy, \methodname is further integrated with a hierarchy memory
framework for personalized presentation generation. The framework has two
levels: long-term memory and working memory. Working memory maintains the
current session state and temporary feedback across revision rounds, so later
local edits can preserve active temporary memory from the same deck.
Long-term memory persists across jobs and is further divided into user profile
memory and tool memory. User profile memory is user-specific and intent-aware:
for each user and each intent, preferences are organized by various
dimensions including theme, content, visual, layout, template, and general, and are routed into working memory when each job starts. 
Tool memory captures reusable execution experience for later edits. This
hierarchy memory framework enables generating and revising presentations according to both
persistent user preferences and the active intent of the current session. The overview of \methodname is shown in Figure~\ref{fig:pptagent_memory_workflow}.

%We evaluate whether \methodname improves personalized presentation generation, rather than only aggregate deck quality or instruction satisfaction \citep{zheng2025pptagent,zheng2026deeppresenter}. 
In experiments, we evaluate \methodname 
across multiple dimensions including slide deck quality, instruction satisfaction, and preference alignment. 
We develop \emph{persona-alignment judgments} as metrics to evaluate personalization alignment. 
%and the \emph{matched local revision comparison} metric for multi-turn modification. 
Qualitative comparisons further illustrate the effectiveness of MemSlides in aligning with users' preferences.
%working-memory preference carryover.

Our contributions are threefold:
\begingroup
\setlength{\parskip}{0pt}
\setlength{\topsep}{2pt}
\setlength{\partopsep}{0pt}
\setlength{\itemsep}{1pt}
\setlength{\parsep}{0pt}
\begin{itemize}
    \item We introduce \methodname, a personalized presentation agent that
    supports multi-turn localized revision. By maintaining session state and
    applying targeted slide-level updates instead of repeated full-deck
    regeneration, \methodname provides the interaction substrate needed for
    learning user preferences from revision feedback.
    \item We further develop a hierarchical memory framework for \methodname consisting of
    a long-term memory and a working memory. 
    The long-term memory contains user profile memory and tool memory, allowing
    stable user preferences and reusable execution experience to persist across
    jobs, while working memory tracks session-specific constraints.
    \item We construct a multi-persona, multi-intent
    user profile bank for personalized presentation generation evaluation.
    Experiments show that user profile memory improves round-0 persona
    alignment, tool memory enhances localized modify reliability in diagnostic
    matched-pair comparisons, and working memory supports
    session-level preference carryover.
\end{itemize}
\endgroup
